Una varilla homogénea de20cm está inicialmente a100°C. En t=0 sus extremos se sumergen
en un baño a0°C y se mantienen allí; la superficie lateral está aislada.
Utilizar el modelo unidimensional con difusividad constante k.
Para plata, aluminio y hierro fundido tomar respectivamente
\(k=1.71,0.86,0.12\ \mathrm{cm^2/s}\).

\begin{enumerate}[label=(\alph*),ref=\alph*,leftmargin=*]
\item Determinar la temperatura en cualquier punto de la varilla para t>0.
\item Usar los dos primeros términos no nulos para estimar u(10,30) en plata, aluminio y hierro fundido.
\end{enumerate}
